\section{Transformers as General-Purpose Sequence Architectures}

The previous section introduced attention as a mechanism for content-based weighted aggregation.
A transformer turns that mechanism into a full architectural template.
Instead of treating attention as an auxiliary device attached to a recurrent model, the transformer makes attention the central structured operator of the network.

\medskip

This shift is conceptually important.
A transformer layer is not merely a collection of engineering tricks.
It is a mathematically analyzable operator on a sequence of representations:
it mixes information across positions by learned pairwise interactions, applies nonlinear pointwise transformations, and stabilizes depth through residual connections and normalization.
Seen this way, transformers are general-purpose sequence architectures because they provide a reusable algebraic block that can model local, global, and content-dependent interaction patterns within the same framework.

\subsection*{Scaled dot-product self-attention}

Let
\[
X \in \mathbb{R}^{T \times d}
\]
be a matrix of token representations, where \(T\) is the sequence length and \(d\) is the representation dimension.
A self-attention layer constructs queries, keys, and values from the same input:
\[
Q = XW_Q, \qquad K = XW_K, \qquad V = XW_V,
\]
where
\[
W_Q \in \mathbb{R}^{d \times d_k},
\qquad
W_K \in \mathbb{R}^{d \times d_k},
\qquad
W_V \in \mathbb{R}^{d \times d_v}.
\]

\begin{definition}[Scaled dot-product attention]
For matrices \(Q \in \mathbb{R}^{T \times d_k}\), \(K \in \mathbb{R}^{T \times d_k}\), and \(V \in \mathbb{R}^{T \times d_v}\), the scaled dot-product attention operator is
\[
\mathrm{Attn}(Q,K,V)
=
\mathrm{softmax}_{\mathrm{row}}\!\left(\frac{QK^{\top}}{\sqrt{d_k}}\right)V,
\]
where \(\mathrm{softmax}_{\mathrm{row}}\) applies the softmax map to each row of the score matrix.
\end{definition}

The matrix
\[
S = \frac{QK^{\top}}{\sqrt{d_k}} \in \mathbb{R}^{T \times T}
\]
contains pairwise similarity scores between positions.
After row-wise softmax one obtains an attention matrix
\[
A = \mathrm{softmax}_{\mathrm{row}}(S),
\]
whose entries satisfy \(A_{ij} \ge 0\) and \(\sum_{j=1}^T A_{ij} = 1\) for every row \(i\).
Thus the output at position \(i\) is
\[
(AV)_i = \sum_{j=1}^T A_{ij} v_j,
\]
which is a convex combination of the value vectors.
This makes self-attention a structured interaction operator:
each output token is built by aggregating information from all positions, but with weights determined adaptively by content.

\medskip

The scaling by \(\sqrt{d_k}\) is not cosmetic.
Without it, dot products can grow in magnitude as the key dimension increases, pushing the softmax into very sharp regimes and making optimization less stable.
The scaling keeps score magnitudes in a more controlled range.

\subsection*{Self-attention as a sequence operator}

A recurrent model updates positions through a temporal chain.
A self-attention layer instead forms an interaction graph in one step.
Every position can communicate directly with every other position through the matrix \(A\).
This is one reason transformers are so expressive for long-range dependence:
if token \(i\) needs information from token \(j\), the model does not have to pass that information through \(j-i\) intermediate recurrent transitions.
It can create a direct content-based edge immediately.

\medskip

In pure self-attention, however, the interaction pattern depends only on the token representations supplied to the layer.
If no positional information is added, the layer does not intrinsically know whether one token came earlier or later in the sequence.
This observation can be stated precisely.

\begin{proposition}[Permutation equivariance of self-attention without positional information]
Let \(P \in \mathbb{R}^{T \times T}\) be a permutation matrix, and define
\[
\mathrm{SA}(X)
=
\mathrm{Attn}(XW_Q, XW_K, XW_V).
\]
If no positional information is added to \(X\), then
\[
\mathrm{SA}(PX) = P\,\mathrm{SA}(X).
\]
That is, self-attention is permutation equivariant.
\end{proposition}

\paragraph{Proof sketch.}
If \(X' = PX\), then
\[
Q' = X'W_Q = PQ,
\qquad
K' = X'W_K = PK,
\qquad
V' = X'W_V = PV.
\]
Therefore the score matrix becomes
\[
S' = \frac{Q'{K'}^{\top}}{\sqrt{d_k}}
= \frac{(PQ)(PK)^{\top}}{\sqrt{d_k}}
= \frac{P Q K^{\top} P^{\top}}{\sqrt{d_k}}
= P S P^{\top}.
\]
Applying row-wise softmax simply permutes rows and columns in the same way, so
\[
\mathrm{softmax}_{\mathrm{row}}(S') = P\,\mathrm{softmax}_{\mathrm{row}}(S)\,P^{\top}.
\]
Hence
\[
\begin{aligned}
\mathrm{SA}(PX)
&=
\mathrm{softmax}_{\mathrm{row}}(S')V' \\
&=
P\,\mathrm{softmax}_{\mathrm{row}}(S)P^{\top}PV \\
&=
P\,\mathrm{softmax}_{\mathrm{row}}(S)V \\
&=
P\,\mathrm{SA}(X).
\end{aligned}
\]
So permuting the inputs permutes the outputs in exactly the same way.
\qed

This proposition is pedagogically valuable because it identifies a precise limitation.
Self-attention alone is naturally suited not only to sequences but also to sets, where permutation equivariance may actually be desirable.
For sequence modeling, however, order matters.
Transformers therefore inject positional information explicitly, either by adding positional encodings or by using learned positional embeddings.
Once order information is included, the architecture is no longer equivariant to arbitrary permutations of the token positions.

\subsection*{Multi-head attention}

A single attention map produces one interaction pattern.
In practice, transformers use \emph{multi-head attention}, which computes several attention operations in parallel with different learned projections.
For head \(h = 1,\dots,H\), define
\[
Q_h = XW_Q^{(h)},
\qquad
K_h = XW_K^{(h)},
\qquad
V_h = XW_V^{(h)},
\]
and then
\[
\mathrm{head}_h(X)
=
\mathrm{Attn}(Q_h,K_h,V_h).
\]
The head outputs are concatenated and linearly mixed:
\[
\mathrm{MHA}(X)
=
\mathrm{Concat}(\mathrm{head}_1(X),\dots,\mathrm{head}_H(X))W_O.
\]

The point is not merely to enlarge the parameter count.
Different heads can represent different types of interactions:
local versus long-range dependence,
syntactic versus semantic relation,
or one-to-one copying versus broad contextual smoothing.
The transformer layer therefore learns not one interaction graph, but a small family of interaction graphs that are recombined into a new representation.

\subsection*{The transformer block written algebraically}

A transformer layer combines self-attention with a positionwise nonlinear map.
A convenient modern form is the \emph{pre-normalization} block.
If \(H^{(\ell)} \in \mathbb{R}^{T \times d}\) is the input to layer \(\ell\), then one writes
\[
U^{(\ell)} = \mathrm{LN}(H^{(\ell)}),
\]
\[
\widetilde{H}^{(\ell)}
=
H^{(\ell)} + \mathrm{MHA}(U^{(\ell)}),
\]
\[
Z^{(\ell)} = \mathrm{LN}(\widetilde{H}^{(\ell)}),
\]
\[
H^{(\ell+1)}
=
\widetilde{H}^{(\ell)} + \mathrm{MLP}(Z^{(\ell)}).
\]
Here the positionwise feedforward network typically has the form
\[
\mathrm{MLP}(z) = W_2\,\sigma(W_1 z + b_1) + b_2,
\]
applied independently to each token representation.

\medskip

This algebra displays the essential architecture clearly.
There are two kinds of transformation in each layer:
\begin{itemize}
    \item \textbf{interaction across positions}, produced by multi-head self-attention,
    \item \textbf{nonlinear feature transformation within each position}, produced by the MLP.
\end{itemize}
Residual connections preserve an identity pathway through the block, which helps optimization in deep networks, while normalization controls the scale and geometry of intermediate activations.

\medskip

The original transformer formulation used a closely related \emph{post-normalization} arrangement.
The difference is architectural rather than conceptual:
both forms combine attention, nonlinear pointwise transformation, skip connections, and normalization into a repeated block.

\begin{figure}[h]
\centering
\fbox{%
\begin{minipage}{0.88\textwidth}
\centering
\textbf{Transformer block (pre-normalization form)}

\vspace{0.5em}
\begin{tabular}{c}
Input \(H^{(\ell)}\) \\
\(\downarrow\) \\
LayerNorm \\
\(\downarrow\) \\
Multi-Head Self-Attention \\
\(\downarrow\) \\
Add residual from \(H^{(\ell)}\) \\
\(\downarrow\) \\
LayerNorm \\
\(\downarrow\) \\
Positionwise MLP \\
\(\downarrow\) \\
Add residual \\
\(\downarrow\) \\
Output \(H^{(\ell+1)}\)
\end{tabular}
\end{minipage}}
\caption{A monochrome block diagram of a transformer layer.
Each block mixes global interaction through self-attention with per-position nonlinear processing through an MLP, and both sublayers are wrapped in residual pathways.}
\end{figure}

\subsection*{A tiny worked example}

A small numerical example helps make the operator concrete.
Consider a sequence of length \(T=3\) with one-dimensional queries and keys:
\[
Q = K =
\begin{bmatrix}
1\\
0\\
1
\end{bmatrix},
\qquad
V = I_3.
\]
For simplicity, take \(d_k = 1\), so no nontrivial scaling appears.
Then
\[
S = QK^{\top}
=
\begin{bmatrix}
1 & 0 & 1\\
0 & 0 & 0\\
1 & 0 & 1
\end{bmatrix}.
\]
Applying row-wise softmax gives
\[
A
=
\mathrm{softmax}_{\mathrm{row}}(S)
\approx
\begin{bmatrix}
0.422 & 0.155 & 0.422\\
0.333 & 0.333 & 0.333\\
0.422 & 0.155 & 0.422
\end{bmatrix}.
\]
Since \(V = I_3\), the attention output is simply
\[
AV = A.
\]
Thus the first and third tokens place most of their weight on positions \(1\) and \(3\), while the second token averages all positions equally.

\begin{example}[Reading the tiny attention matrix]
This matrix already illustrates several essential facts.
First, positions \(1\) and \(3\) can interact directly in one layer even though they are not adjacent.
Second, the rows of \(A\) sum to one, so each output is a weighted average of values.
Third, without positional information, the first and third positions behave identically because their representations are identical.
A positional encoding would break this symmetry and allow the model to distinguish ``first'' from ``third.'' 
\end{example}

One should not mistake this toy example for a realistic language model computation.
Its purpose is structural.
It shows that self-attention is an explicit matrix operator whose pattern of interaction can be inspected directly.
In larger models the same principle survives, only with higher-dimensional features and many heads.

\subsection*{Complexity and comparison with recurrence}

Transformers are often contrasted with recurrent models because both are used for sequential data, but the contrast is subtler than ``new versus old.''
They organize computation differently.

\medskip

For a recurrent layer with hidden width \(d\), the dependence on sequence length is typically linear:
processing \(T\) positions requires \(T\) recurrent updates, each using matrix multiplications of size roughly \(d \times d\).
This leads to arithmetic cost on the order of \(O(Td^2)\), with an inherently sequential computation graph.

\medskip

For dense self-attention, one still pays \(O(Td^2)\) for the linear projections, but the interaction matrix \(QK^{\top}\) and the multiplication by \(V\) contribute an additional \(O(T^2 d)\) term.
Thus, with width fixed, dense self-attention has quadratic dependence on sequence length.
The price of direct all-to-all interaction is the \(T^2\) attention matrix.

\medskip

Yet there is an architectural advantage.
The pairwise interactions for all positions can be computed in parallel once the token representations for a layer are available.
Moreover, the path length between two arbitrary positions is very short:
one self-attention layer can connect any token directly to any other token, whereas a recurrent model may require many sequential steps for the same information to propagate.

\begin{center}
\begin{tabular}{|p{0.23\textwidth}|p{0.22\textwidth}|p{0.18\textwidth}|p{0.22\textwidth}|}
\hline
Architecture & Cost in \(T\) (with width fixed) & Sequential depth & Long-range path length \\
\hline
Recurrent layer & linear in \(T\) & \(T\) & on the order of \(T\) \\
Dense self-attention layer & quadratic in \(T\) & essentially constant per layer & \(1\) \\
\hline
\end{tabular}
\end{center}

This comparison explains both the success and the limitations of transformers.
They are powerful because they make global interaction easy and parallelizable.
They can be expensive because unrestricted interaction across all token pairs scales poorly for very long sequences.
Many later architectural refinements can be understood as attempts to preserve the expressive advantages of attention while reducing this quadratic cost.

\subsection*{Why transformers became general-purpose}

The transformer block is now used far beyond classical text sequences.
It can process tokenized image patches, audio frames, multimodal streams, sets, and structured collections of objects.
The reason is not that all such data are ``really language.''
The reason is that the block itself is abstract and modular.
It only assumes that we have a collection of representations together with some way of encoding whatever structural information matters --- positional, spatial, temporal, or relational.

\medskip

From this point of view, a transformer is a general-purpose architecture because it separates three ingredients:
\begin{itemize}
    \item a representation of the elements to be processed,
    \item an interaction operator that mixes information across elements by learned relevance,
    \item a per-element nonlinear map that enriches features after interaction.
\end{itemize}
That decomposition is flexible enough to be reused across many domains.

\paragraph{What to remember.}
A transformer layer is a repeated algebraic block built from multi-head self-attention, a positionwise MLP, residual connections, and normalization.
Scaled dot-product attention computes a row-stochastic interaction matrix from learned similarities.
Without positional information, self-attention is permutation equivariant, so explicit positional encoding is necessary for genuine sequence modeling.
Transformers replace long sequential interaction paths with direct pairwise interaction, but dense self-attention incurs quadratic cost in sequence length.

\paragraph{Common confusion.}
It is easy to say that transformers ``remove order.''
That is not quite right.
The attention operator by itself is insensitive to order, but the full model can represent order once positional information is supplied.
It is also easy to think that the transformer is ``just attention.''
In fact, the architecture depends equally on residual structure, normalization, and the repeated alternation between cross-position interaction and per-position nonlinear transformation.

\paragraph{Exercises.}
\begin{enumerate}
    \item Let \(X \in \mathbb{R}^{T \times d}\) and let \(P\) be a permutation matrix.
    Write out carefully why \(QK^{\top}\) transforms as \(P(QK^{\top})P^{\top}\) under the substitution \(X \mapsto PX\).
    \item Suppose all token representations in a self-attention layer are identical and no positional information is added.
    Show that every row of the attention matrix is the same.
    What does this imply about the outputs?
    \item For the toy example in this section, replace
    \[
    Q = K = [1,0,1]^{\top}
    \]
    by
    \[
    Q = K = [2,0,-2]^{\top}.
    \]
    Compute the new attention matrix and explain how the sharper score differences change the interaction pattern.
    \item Compare a recurrent layer and a self-attention layer as mechanisms for modeling a dependency between the first and last token of a long sequence.
    Which architecture has the shorter interaction path, and which has the cheaper asymptotic dependence on sequence length?
    \item Derive the output dimensions of a multi-head attention block with \(H\) heads, model width \(d\), per-head key dimension \(d_k\), and per-head value dimension \(d_v\).
\end{enumerate}

\paragraph{Notes and references.}
The transformer architecture was introduced by Vaswani and collaborators in \emph{Attention Is All You Need}.
Its attention mechanism builds on earlier encoder--decoder attention ideas, especially the work of Bahdanau, Cho, and Bengio on neural machine translation.
For pedagogical development, it is useful to view the transformer not as a historical slogan but as a structured operator composed of attention, residual pathways, normalization, and per-position nonlinear maps.
That viewpoint also clarifies why positional encoding is essential and why complexity in sequence length became a central issue in later transformer variants.
